\documentclass{article}
\usepackage{amssymb}
\usepackage{amsmath}
\usepackage{geometry}
\usepackage{parskip}
\usepackage{amsthm}

\usepackage[style=ieee]{biblatex}
\addbibresource{references.bib}

\geometry{
 a4paper,
 total={170mm,257mm},
 left=20mm,
 top=20mm,
 }

\newcommand{\Q}{\mathbb{Q}}
\newcommand{\Z}{\mathbb{Z}}
\newcommand{\N}{\mathbb{N}}
\newcommand{\abs}[1]{\left| #1 \right|}

\title{Arithmetic Geometry - Sheet 02}
\author{Andres Felipe Jaramillo Garces}
\date{November 2026}

\begin{document}
\maketitle
\section*{Exercise 1}
\begin{proof}

We start by noting that $\frac{\log\abs{\alpha - \frac{a}{q}}}{\log q} = \log_q\abs{\alpha - \frac{a}{q}}$. If this is 1, then this is the same as saying $1/q = \abs{\alpha - \frac{a}{q}}$. In the language of the proof of Dirichlet's theorem that we saw in class, this is the same as saying that we can get "within a box's distance" of $\alpha$ using the "boxes" of $a/q$. In this sense, a bigger $q$ implies smaller, more precise "boxes" that could, intuitively, get closer to $\alpha$, as they divide the number line in smaller portions.

$\rightarrow$ Let $\alpha = \frac{c}{d} \in \Q, a, q \in \Z$.
It is clear that $\abs{\alpha - \frac{a}{q}} \in \Q$. Further, by solving the inequality $\abs{\frac{c}{d} - \frac{a}{q}} \leq 1/q$ for $a$ we obtain that $\frac{qc}{d} - 1 \leq a \leq \frac{qc}{d} + 1$, which is always solveable for some integer as long as $q$ is big enough. Thus, for $q$ sufficiently big, it it always possible to get "within a $q$-box's distance" of $\alpha$, implying $\forall_{q \in \N} \exists_{a \in \Z} \frac{\log\abs{\alpha - \frac{a}{q}}}{\log q} \leq -1$.

\end{proof}


\section*{Exercise 3}
\textbf{Note:} I haven't been to class this week and I have no idea what has been taught or not. After looking around for a while, I found \cite{Everste1}, of which Prof. Viada uploaded the first part. I then just read Chapter 3, sections 3.1, 3.2, and 3.5. Turns out a more general case of this problem is one of the exercises in section 3.7 there. I use this reference throughout this proof. I made myself an indexed PDF of the lecture notes that I can gladly share with the class.


\begin{proof}
Since $f \in \Z[x]$, we conclude that $\alpha_1$ and $\alpha_2$ are algebraic integers (See \cite[41]{Everste1}), meaning there is a polynomial with coefficients in $\Z$ for which they are zeros. Consider now $\beta = \alpha_i - p/q$ with $p/q \in \Q$ and multiply by $q$ to produce $q\beta = \alpha_i - p$.  Both $\alpha_i$ and $p$ are algebraic integers (For $p$ consider $g = x - p \in \Z[x]$). As per \cite[\nopp Lemma 3.3]{Everste1}, we obtain that $q\beta$ is an algebraic integer, since it turns out algebraic integers are a closed ring. 

Because $\beta \in \Q[\alpha_1, \alpha_2]$, we know that the minimal polynomial of $\beta$ (call it $h_{\beta}$) has a degree at most as high as the extension, and that is degree 2 (See \cite[\nopp Corollary 3.23]{Everste1}). Observe that $f(x + p/q)$ is a polynomial of degree 2 for which $\beta$ is a zero. Thus $h_{\beta} = f(x + p/q)$. By constructing $h(x) = h_{\beta}(x/q)$, we obtain a polynomial of degree 2 of which $q\beta$ is a 0. Since clearly $q\beta \notin \Q$, the degree of its minimal polynomial must be at least 2. Thus, $h$ is the minimal polynomial of $q\beta$. Further, as per \cite[\nopp Lemma 3.4]{Everste1}, $h \in \Z[x]$. 

Let $\gamma$ be the other root of $h_{\beta}$. Clearly $q\gamma$ is also a zero of $h$. Thus we conclude that $h(0) = q^2 \gamma \beta \in \Z$ an in particular $\abs{q^2 \gamma \beta} \geq 1$. Define $\overline{|\beta|} = \max(\abs{\gamma}, \abs{\beta})$. We obtain

\begin{align*}
  \abs{q^2 \gamma \beta} &\geq 1 \\
  q^2 \overline{|\beta|} \abs{\beta} &\geq 1 \\
  \abs{\beta} &\geq q^{-2} \gamma^{-1} \\
  \abs{\alpha_i - p/q} &\geq q^{-2} \gamma^{-1} \\
\end{align*}

We're not done becasue $\gamma$ still depends on $p/q$. Consider that for this problem to be interesting, we're looking for pairs $p/q$ for which $\abs{\alpha_i - p/q} < 1$, since otherwise the inequality is true for any constant smaller than $q^2$ (pick $C = 1$). Solving this inequality for $p/q$ gives us $\alpha_i - 1 \leq p/q \leq \alpha_i + 1$. Going back to the way that we constructed $h_{\beta}$, we notice that this defines a 'bounded' (by size of coefficients) set of possible polynomials for $\beta$. We simply need to pick the minimal corresponding $\gamma$ from the polynomials of this set and set $C$ to that.
\end{proof}
\printbibliography

\end{document}

